\def\writeANS{}
\begin{loigiaibth}{31}
  $80\,\mathrm {g}$ $\mathrm {NH_4NO_3} \sim 1\,\mathrm {mol} \Rightarrow Q = 26\,\mathrm {(kJ)}$ \\ $\Delta H > 0$, quá trình hòa tan thu nhiệt, nhiệt độ giảm đi một lượng là: \\ $\Delta T = \dfrac {26.10^3}{4{,}2.10^3} = 6{,}2^\circ \mathrm {C} \Rightarrow $ Nhiệt độ cuối cùng là $25 - 6{,}2 = 18{,}8^\circ \mathrm {C}$
\end{loigiaibth}
\def\writeANS{}
\begin{loigiaibth}{32}
  Phản ứng xảy ra: $\mathrm {HCl(aq) + NaHCO_3(aq) \longrightarrow NaCl\,(aq) + H_2O(l) + CO_2(g)}$ \\ $\Delta H = (-407) + (-286) + (-392) - (-168) - (-932) = 15\,\mathrm {(kJ)}$ \\ $\Rightarrow $ Phản ứng thu nhiệt. \\ Số mol $\mathrm {HCl} =$ số mol $\mathrm {NaHCO_3} = 0{,}1\,\mathrm {mol} \Rightarrow Q = 0{,}1.15 = 1{,}5\,\mathrm {(kJ)}$ \\ Nhiệt độ giảm đi: $\Delta T = \dfrac {1{,}5.10^3}{200.4{,}2} = 1{,}8^\circ \mathrm {C}$ \\ Nhiệt độ ban đầu khi trộn hai dung dịch: $\dfrac {27 + 28}{2} = 27{,}5^\circ \mathrm {C}$ \\ $\Rightarrow $ Nhiệt độ cuối cùng là: $27{,}5 - 1{,}8 = 25{,}7^\circ \mathrm {C}$
\end{loigiaibth}
\def\writeANS{}
\begin{loigiaibth}{33}
  Khi trộn hai dung dịch, nhiệt độ trước phản ứng là: $\dfrac {25+26}{2} = 25{,}5^\circ \mathrm {C}$ \\ Nhiệt lượng tỏa ra là: $Q = (50 + 50).4{,}2.(28 – 25{,}5) = 1050\,\mathrm {(J)}$ \\ Phản ứng xảy ra: $\mathrm {AgNO_3(aq) + NaCl(aq) \longrightarrow AgCl(s) + NaNO_3(aq)}$ \\ Số mol $\mathrm {AgNO_3} =$ số mol $\mathrm {NaCl} = \dfrac {0{,}5.50}{1000} = 0{,}025\,\mathrm {mol}$ \\ $\Rightarrow \Delta H = \dfrac {1050}{0{,}025} = 42000\,\mathrm {J} = 42\,\mathrm {(kJ)}$
\end{loigiaibth}
